\subsection{Packages và Pages: gói NuGet nội bộ và trang tài liệu}

GitHub không chỉ lưu mã nguồn mà còn cung cấp kho gói tích hợp mang tên
GitHub Packages, cho phép nhóm đồ án publish các thư viện dùng chung ở dạng
gói \verb|*.nupkg| thay vì chép file \verb|.dll| thủ công qua USB hay email.
Mỗi gói gắn với một repository, kế thừa quyền truy cập của repository đó,
nên thành viên nào đọc được code thì mới tải được gói nội bộ \cite{githubdocs}.

Quy trình publish gói cho project ASP.NET Web Forms gồm ba bước chính:
đóng gói project thư viện bằng \verb|dotnet pack| hoặc \verb|nuget pack|,
thêm nguồn gói trỏ về \verb|nuget.pkg.github.com|, rồi đẩy gói lên bằng
\verb|dotnet nuget push| với token xác thực \verb|GITHUB_TOKEN|.
Token này cần quyền \verb|write:packages|, còn khi cài gói chỉ cần
quyền \verb|read:packages| nên rất phù hợp để cấp cho máy cá nhân của từng
thành viên mà không lo lộ quyền ghi \cite{githubdocs}.

\begin{lstlisting}[language=bash]
# Buoc 1: dong goi thu vien dung chung (ASCII only comments)
dotnet pack MyLib/MyLib.csproj -c Release -o ./nupkg
# Buoc 2: them nguon GitHub Packages (OWNER = ten user hoac org)
dotnet nuget add source https://nuget.pkg.github.com/OWNER/index.json \
  --name github --username OWNER --password $GITHUB_TOKEN --store-password-in-clear-text
# Buoc 3: day goi len GitHub Packages
dotnet nuget push ./nupkg/MyLib.1.0.0.nupkg --source github
# Cai goi tu may thanh vien khac trong nhom
dotnet add package MyLib --version 1.0.0 --source https://nuget.pkg.github.com/OWNER/index.json
\end{lstlisting}

Với project Web Forms trên .NET Framework 4.8.1, lệnh \verb|dotnet pack|
không dùng trực tiếp cho project \verb|.aspx| mà dùng cho project thư viện
\verb|Class Library| đi kèm, ví dụ project chứa lớp tiện ích mã hóa mật khẩu
hay lớp phân trang dùng chung cho ba thành viên \cite{msnetfx}.
Phiên bản gói nên đặt theo quy tắc Semantic Versioning, ví dụ \verb|1.2.0|,
kết hợp Conventional Commits để sinh changelog tự động \cite{convcommits}.

GitHub Pages là dịch vụ host trang web tĩnh miễn phí ngay từ repository,
rất phù hợp để công bố tài liệu đồ án: đặc tả yêu cầu, sơ đồ cơ sở dữ liệu,
hướng dẫn cài đặt và nhật ký demo từng tuần.
Nhóm chỉ cần đưa file HTML tĩnh hoặc source Jekyll vào nhánh \verb|main|
hoặc thư mục \verb|docs|, bật Pages trong phần Settings là có ngay địa chỉ
\verb|https://minhtritm.github.io/AspxDemo_NET481/| để nộp kèm báo cáo.

Cách hiện đại nhất là deploy Pages bằng GitHub Actions: mỗi lần merge vào
nhánh \verb|main|, workflow tự động build tài liệu và publish, khỏi cần giữ
nhánh \verb|gh-pages| thủ công như trước đây \cite{ghactions}.
Ví dụ tối thiểu dưới đây build thư mục \verb|docs| và đẩy lên Pages:

\begin{lstlisting}[language=yaml]
# Workflow deploy tai lieu do an len GitHub Pages
name: docs-to-pages
on:
  push:
    branches: [main]
    paths: ['docs/**']
permissions:
  contents: read
  pages: write
  id-token: write
jobs:
  deploy:
    runs-on: ubuntu-latest
    environment:
      name: github-pages
    steps:
      - uses: actions/checkout@v4
      - uses: actions/configure-pages@v5
      - uses: actions/upload-pages-artifact@v3
        with:
          path: docs
      - uses: actions/deploy-pages@v4
\end{lstlisting}

Tóm lại, Packages giải quyết bài toán chia sẻ thư viện nội bộ có phiên bản,
còn Pages giải quyết bài toán công bố tài liệu công khai với chi phí bằng không.
Cả hai đều tận dụng sẵn quyền truy cập của repository nên nhóm không cần
dựng thêm server riêng trong phạm vi đồ án môn học.

\subsection{Codespaces: môi trường phát triển trên mây}

GitHub Codespaces là môi trường phát triển đầy đủ chạy trên mây, mở trực tiếp
từ repository bằng trình duyệt hoặc Visual Studio Code \cite{githubdocs}.
Mỗi codespace là một container có sẵn Git, .NET SDK, Node và các extension,
được mô tả bằng file \verb|devcontainer.json| nên mọi thành viên mở lên đều
có môi trường giống hệt nhau, chấm dứt cảnh máy chạy được máy thì không.

Với đồ án ASP.NET Web Forms 4.8.1, cần lưu ý Codespaces chạy trên Linux nên
không build trực tiếp project Web Forms vốn cần MSBuild trên Windows.
Nhóm vẫn dùng Codespaces hiệu quả cho các việc không cần build Windows:
viết báo cáo LaTeX, soạn tài liệu, review pull request, viết YAML cho Actions
và chạy các script kiểm tra chính tả hay kiểm tra link hỏng \cite{msmsbuild}.
Việc build thật project \verb|.aspx| vẫn giao cho GitHub Actions với runner
\verb|windows-latest|, như vậy mỗi người đóng góp đúng thế mạnh của từng nền tảng.

Ba thao tác Codespaces dùng nhiều nhất trong đồ án:

\begin{itemize}
\item Tạo codespace mới từ nhánh \verb|main| hoặc từ một pull request để review code.
\item Mở terminal trong codespace, chạy \verb|git status| và \verb|git log| như máy local.
\item Commit và push về repository, pipeline Actions tự động chạy kiểm tra.
\end{itemize}

\begin{lstlisting}[language=bash]
# Cai dat CLI de dieu khien Codespaces tu terminal
gh extension install github/gh-codespace
# Liet ke cac codespace hien co
gh codespace list
# Tao codespace moi tu repo do an (nhanh main)
gh codespace create --repo OWNER/AspxDemo_NET481 --branch main
# Mo SSH vao codespace de lam viec
gh codespace ssh --codespace my-codespace-name
# Xoa codespace khi xong viec de tiet kiem quota
gh codespace delete --codespace my-codespace-name
\end{lstlisting}

\begin{figure}[H]
\centering
\begin{tikzpicture}[>=Stealth, thick, box/.style={draw, rounded corners, fill=blue!10, minimum width=2.4cm, minimum height=0.9cm, align=center, font=\small}]
\node[box] (a) {Codespaces};
\node[box, right=1.2cm of a] (b) {git push};
\node[box, right=1.2cm of b] (c) {Actions};
\draw[->] (a) -- node[midway, above, font=\footnotesize] {commit} (b);
\draw[->] (b) -- node[midway, above, font=\footnotesize] {build} (c);
\end{tikzpicture}
\caption{Luong lam viec tu Codespaces den Actions}
\label{fig:c43-codespaces-flow}
\end{figure}

Hình trên tóm tắt luồng làm việc: soạn thảo trên Codespaces, đẩy code về
repository, Actions trên runner Windows build project ASPX 4.8.1 và báo kết quả.
Sinh viên dùng gói GitHub Education được cấp quota Codespaces miễn phí mỗi tháng,
đủ cho nhu cầu viết tài liệu và review code của đồ án (xem mục Giáo dục ở cuối).

\subsection{Copilot: trợ lý gợi ý code bằng AI}

Muse là trợ lý lập trình dùng mô hình ngôn ngữ lớn,
gợi ý ngay trong editor khi sinh viên gõ comment hoặc tên hàm \cite{githubdocs}.
Với đồ án Web Forms, Copilot giúp sinh code lặp lại nhanh hơn khoảng 30\%,
ví dụ tạo validator, sinh truy vấn, viết unit test cho lớp tiện ích.
Tuy nhiên Copilot chỉ gợi ý chứ không hiểu nghiệp vụ, nên mọi đoạn code sinh ra
đều phải đọc lại, chạy thử và review như code của người viết.

Ví dụ: sinh viên gõ comment mô tả hàm kiểm tra đăng nhập, Copilot gợi ý khung
hàm C\# kiểm tra tên đăng nhập và mật khẩu cho trang \verb|Login.aspx|.
Đoạn code dưới đây mô phỏng kết quả gợi ý sau khi đã chỉnh sửa để an toàn hơn:

\begin{lstlisting}[language=[Sharp]C]
// Check login for page Login.aspx (demo only)
public bool KiemTraDangNhap(string tenDangNhap, string matKhau)
{
    if (string.IsNullOrWhiteSpace(tenDangNhap)) return false;
    if (string.IsNullOrEmpty(matKhau)) return false;
    string sql = "SELECT MatKhauHash, Muoi FROM NguoiDung WHERE TenDangNhap = @u";
    using (var conn = new SqlConnection(_chuoiKetNoi))
    using (var cmd = new SqlCommand(sql, conn))
    {
        cmd.Parameters.AddWithValue("@u", tenDangNhap);
        conn.Open();
        using (var rd = cmd.ExecuteReader())
        {
            if (!rd.Read()) return false;
            string hashLuu = rd.GetString(0);
            string muoi = rd.GetString(1);
            string hashNhap = HamBamMatKhau(matKhau, muoi);
            return hashLuu == hashNhap;
        }
    }
}
\end{lstlisting}

Ba điểm cần tự kiểm tra sau khi nhận gợi ý từ Copilot:

\begin{enumerate}
\item Truy vấn phải tham số hóa bằng \verb|@u| như trên, tuyệt đối không nối chuỗi
tên đăng nhập vào câu lệnh SQL vì sẽ dính lỗi SQL injection.
\item Mật khẩu phải băm kèm muối bằng thuật toán mạnh, không lưu hay so sánh
mật khẩu gốc, không gắn cứng mật khẩu mẫu trong code.
\item Chuỗi kết nối lấy từ file cấu hình hoặc Secret, không dán trực tiếp vào
code hay vào prompt chat với Copilot.
\end{enumerate}

Về bản quyền và bảo mật, nhóm cần nhớ ba nguyên tắc \cite{githubdocs}:

\begin{itemize}
\item Không dán mã nguồn độc quyền, đề thi hay dữ liệu cá nhân vào prompt,
vì nội dung prompt có thể được dùng theo chính sách của nhà cung cấp.
\item Đoạn gợi ý dài giống hệt code nổi tiếng cần kiểm tra giấy phép trước khi
dùng, tránh vi phạm bản quyền của project nguồn mở.
\item Secret như token, connection string, khóa API không bao giờ xuất hiện
trong prompt, trong code mẫu hay trong lịch sử commit.
\end{itemize}

Khi dùng đúng cách, Copilot là gia sư kiên nhẫn: giải thích code khó,
gợi ý test case, đề xuất refactor tên biến rõ nghĩa hơn cho sinh viên năm nhất.

\subsection{Bảo mật nâng cao: quét secret, Dependabot và CodeQL}

Repository đồ án thường chứa file \verb|Web.config| với chuỗi kết nối cơ sở dữ liệu,
nên tính năng bảo mật của GitHub có giá trị thực tế ngay từ buổi đầu \cite{githubdocs}.
Secret scanning tự động quét lịch sử commit để tìm token, khóa API hay mật khẩu
bị vô tình đẩy lên, rồi gửi cảnh báo cho người maintain repository.

Mạnh hơn một bước, push protection chặn ngay tại thời điểm push: nếu commit chứa
ví dụ một connection string của SQL Server, lệnh push bị từ chối kèm hướng dẫn
xóa secret khỏi commit, thay vì để lọt lên server rồi mới cảnh báo sau.
Ví dụ dưới đây mô phỏng lần push bị chặn và cách khắc phục:

\begin{lstlisting}[language=bash]
# Vo tinh commit file cau hinh chua secret that
git add Web.config
git commit -m "Update connection string"
git push origin feature/login
# Ket qua: remote tu choi vi phat hien secret (push protection)
# remote: error: GH013: Repository rule violations found for refs/heads/feature/login
# Giai phap: dua secret ra khoi code, doc tu bien moi truong
# 1) Xoa commit chua secret (nhanh ca nhan, chua merge)
git reset --soft HEAD~1
# 2) Thay connection string that bang gia tri gia an toan
# 3) Commit lai va push, sau do xoay (rotate) secret bi lo
git commit -m "fix: move connection string to environment variable"
git push origin feature/login
\end{lstlisting}

Dependabot tự động mở pull request nâng cấp các gói phụ thuộc có lỗ hổng,
ví dụ nâng cấp thư viện JavaScript trong thư mục \verb|Scripts| hay gói NuGet
cũ trong project \verb|packages.config| \cite{githubdocs}.
Nhóm chỉ cần review và merge pull request của Dependabot như pull request thường,
pipeline Actions sẽ build lại để xác nhận bản nâng cấp không phá vỡ đồ án.

CodeQL là công cụ phân tích tĩnh: nó dựng cơ sở dữ liệu ngữ nghĩa của mã nguồn
rồi chạy các truy vấn tìm mẫu lỗ hổng như SQL injection hay cross-site scripting.
CodeQL chạy như một workflow Actions, báo kết quả ngay trong thẻ Security của
repository, rất phù hợp để rà soát trang đăng nhập trước buổi demo \cite{ghactions}.

\begin{table}[H]
\centering
\caption{Bon tinh nang bao mat nang cao cua GitHub}
\label{tab:c43-bao-mat}
\begin{tabular}{llll}
\toprule
Tính năng & Phát hiện gì & Chặn ở đâu & Ai xử lý \\
\midrule
Secret scanning & Token và khóa API bị lộ & Sau khi push & Người maintain \\
Push protection & Secret trong commit mới & Trước khi push & Người push \\
Dependabot & Gói phụ thuộc có lỗ hổng & Pull request mới & Người review \\
CodeQL & Mẫu code dễ bị tấn công & Workflow Actions & Cả nhóm \\
\bottomrule
\end{tabular}
\end{table}

Bảng trên giúp nhóm phân công rõ ràng: người push chịu trách nhiệm gỡ secret,
người review chịu trách nhiệm merge bản vá Dependabot, cả nhóm cùng đọc báo cáo
CodeQL trước mỗi mốc nộp bài.

\subsection{Gói GitHub Education cho sinh viên}

GitHub Education là gói ưu đãi miễn phí dành cho sinh viên, giảng viên và nhà
trường, gom nhiều công cụ trả phí thành một gói duy nhất \cite{githubdocs}.
Với sinh viên, ba lợi ích thiết thực nhất cho đồ án ASPX là: Copilot miễn phí
trong thời gian học, quota Codespaces miễn phí mỗi tháng và không giới hạn
repository private cùng cộng tác viên cho nhóm đồ án.

Ngoài ra gói còn kèm nhiều ưu đãi từ đối tác như tên miền miễn phí, credit
cloud để demo và khóa học bổ trợ, giúp nhóm công bố sản phẩm môn học nghiêm túc
mà không tốn chi phí hạ tầng.

Các bước đăng ký gói sinh viên:

\begin{enumerate}
\item Chuẩn bị email do trường cấp, ví dụ email có đuôi \verb|@edu.vn| hoặc
\verb|@st.university.edu.vn|, kèm thẻ sinh viên hoặc giấy xác nhận đang học.
\item Truy cập trang \url{https://education.github.com/} và chọn mục đăng ký
gói Student Developer Pack.
\item Điền thông tin trường, tải minh chứng sinh viên và chờ GitHub xét duyệt,
thường từ vài giờ đến vài ngày.
\item Sau khi được duyệt, bật Copilot và kiểm tra quota Codespaces trong phần
cài đặt tài khoản, rồi mời hai thành viên còn lại cộng tác repository đồ án.
\end{enumerate}

Lưu ý: ưu đãi gắn với trạng thái sinh viên nên cần gia hạn sau mỗi chu kỳ học,
khi tốt nghiệp một số tính năng sẽ chuyển về gói miễn phí thông thường.
Nhóm nên đăng ký sớm từ đầu học kỳ để tận dụng Copilot và Codespaces cho toàn bộ
quá trình làm đồ án thay vì chờ đến tuần nộp bài mới đăng ký \cite{githubdocs}.

Như vậy, từ Packages, Pages, Codespaces, Copilot đến bộ công cụ bảo mật và gói
Education, GitHub cung cấp một hệ sinh thái khép kín cho nhóm sinh viên: viết code,
chia sẻ thư viện, công bố tài liệu, rà soát bảo mật và demo sản phẩm trên một
nền tảng duy nhất, sẵn sàng liên thông sang GitLab ở chương tiếp theo.
